\begin{abstract}
This paper presents a comparative benchmark of dynamic frequency estimators for low-inertia inverter-based-resource (IBR) grids under both standards-style tests and composite inverter-driven disturbances. The study covers \ArchivalScenarioCount{} scenarios, \ArchivalEstimatorCount{} estimators, \ArchivalPairCount{} tuned estimator-scenario configurations, and \ArchivalMonteCarloRuns{} realizations per configuration. Each estimator-scenario pair is tuned separately and then evaluated under a common dual-rate signal chain in which physics-rate waveforms are synthesized at \SI{1}{MHz} and consumed by the estimators on decimated \SI{10}{kHz} inputs. The paper reports six headline metrics spanning error, relay-oriented trip-risk, recovery, and computational cost, and it aggregates results by disturbance family so that ramp tracking, amplitude steps, modulation, phase discontinuities, spectral interference, ringdown, and multi-event stress remain distinguishable in the final interpretation. The main result is not a single leaderboard but a disturbance-dependent latency-robustness-risk map: SOGI-PLL leads IEEE ramps, EKF leads magnitude steps and out-of-band interference, RA-EKF leads modulation, SOGI-FLL leads phase jumps, UKF leads IBR ringdown, and PLL leads the composite IBR multi-event case on family-level RMSE. Trip-risk leaders often differ from RMSE leaders, and the hardest scenarios are the multi-event sequence and the two \SI{60}{\degree} phase-jump cases. These results show that estimator selection in low-inertia grids is governed jointly by disturbance family, protection-oriented risk, and computational budget, and that compliance-style summaries alone are insufficient for deployment decisions.
\end{abstract}
